\documentclass[11pt,a4paper]{article}
\usepackage[utf8]{inputenc}
\usepackage[T1]{fontenc}
\usepackage[spanish]{babel}
\usepackage{hyperref}
\usepackage{geometry}
\geometry{margin=2.5cm}
\usepackage{parskip}

\title{011 devolver las armas}
\author{Alejandro Valenzuela Rivera\\\small Universidad Autónoma de Nuevo León}
\date{\today}
\begin{document}
\maketitle

\begin{quote}
\textbf{Sesgo declarado.} Este trabajo se sostiene desde el sesgo declarado de su autor, y no finge neutralidad. Lo que afirma se puede tumbar, y ese es el punto: el falsificador está escrito, no escondido.
\end{quote}

\section*{Resumen}
La idea se entrega con su metodo de contestacion, no con la vista.

\section{Por qué se sostiene}
La academia tiene la puerta; el mecanismo la deja abierta.

\section{De dónde viene}
All-Star Superman.

\section{El caso}
All-Star Superman, Lex Luthor: la vista que llega tarde.

\section{Método}
Una capacidad que se devuelve cuando el fin termino.

\section{Límites, y qué lo tumba}
Una capacidad que se queda porque ya se tiene.

\section{Lo que queda abierto}
\textquestiondown{}Como se hereda una vista sin preguntar por el autor?

\section*{Cómo se leyó esta tesis}
Esto se escribió en el repositorio man-made-empathy, del que este candado es una pieza: candados/06 del caso/011 devolver las armas.md. La máquina que lo compuso lo hizo desde este mismo archivo, y no agregó nada que no estuviera escrito. Si algo de esto es falso, se dice en el mismo sitio, y se corrige con fecha.

\end{document}
